\section{De la demo a producción: el verdadero umbral del App Builder}

\subsection{Por qué el éxito del prototipo no equivale al éxito del producto}

Uno de los puntos que más se pasa por alto en el curso es que aquello en lo que el App Builder de IA es más hábil es en hacer visible con rapidez una «idea de interfaz»; pero el trabajo real de conversión en producto sigue existiendo en gran medida.

\begin{itemize}
    \item En la etapa de prototipo basta con hacer que la página y la interacción funcionen
    \item La etapa de producto exige permisos, gestión de estado, manejo de errores, monitoreo, pruebas y mantenibilidad
    \item El lugar donde el equipo realmente invierte tiempo no suele ser generar la primera página de la interfaz, sino conectar el resultado generado con los sistemas de negocio reales
\end{itemize}

\begin{warningbox}{Entre más rápida sea la primera versión generada por la IA, más fácil es subestimar el costo de integración posterior}
Muchos equipos, al ver en cuestión de minutos una «página que funciona», caen en la ilusión de que ya se completó el 80\% del trabajo. Pero en un sistema de producción, lo que en realidad determina el costo suele ser el 20\% restante: el contrato de datos, los límites de error, el modelo de permisos, la revisión y la reversión.
\end{warningbox}

\subsection{El sistema de diseño y el sistema de ingeniería deben volver a conectarse}

Para que el App Builder entre de verdad al flujo de trabajo de un equipo, por lo general debe conectar dos sistemas:

\begin{enumerate}
    \item \textbf{El sistema de diseño}: tokens de color, especificaciones de componentes, tipografía, restricciones de interacción
    \item \textbf{El sistema de ingeniería}: versión del framework, estructura de directorios, patrones de obtención de datos, pruebas y despliegue
\end{enumerate}

\begin{table}[H]
\centering
\begin{tabular}{p{0.2\textwidth} p{0.28\textwidth} p{0.28\textwidth}}
\toprule
Capa del sistema & Qué pasa si falta & Efecto sobre el App Builder \\
\midrule
Sistema de diseño & La página funciona pero la apariencia visual no es consistente & Cada generación parece un producto distinto ensamblado \\
Especificación de componentes & El código es difícil de reutilizar y la interacción es inconsistente & La IA no puede reutilizar de forma estable las piezas ya existentes \\
Contrato de datos & La interfaz queda desconectada del backend real & Solo se puede permanecer en la capa de la demo simulada \\
Pruebas y despliegue & El resultado generado no puede publicarse de forma estable & El equipo no se atreve a conectar el resultado de la IA al flujo principal \\
\bottomrule
\end{tabular}
\caption{Capas del sistema que deben completarse para pasar de herramienta de prototipo a herramienta de producción}
\end{table}

\begin{knowledgebox}{El mejor App Builder no sustituye al sistema de ingeniería, sino que acorta la distancia para conectarse a él}
Las herramientas verdaderamente destacadas no fingen que la complejidad de la ingeniería no existe, sino que procuran alinear el resultado generado con el sistema de diseño, el sistema de tipos, el proceso de pruebas y la forma de despliegue que el equipo ya tiene. Así, la IA es un acelerador y no una nueva isla aislada.
\end{knowledgebox}

\subsection{Resumen de la sección}

El valor central del App Builder está en convertir con rapidez una idea en una interfaz que se puede discutir, pero para cambiar de verdad el flujo de trabajo de un equipo debe atravesar umbrales más duros: el sistema de diseño, el contrato de datos, las pruebas y el despliegue.

